\expandafter\ifx\csname CAPETINCLUDEMODE\endcsname\relax
\documentclass[11pt,a4paper]{article}
% preamble.tex — packages et commandes communes a toutes les fiches CAPET SII IE
\input{glyphtounicode}
\pdfgentounicode=1
\usepackage{cmap}
\usepackage[T1]{fontenc}
\usepackage[utf8]{inputenc}
\usepackage[french]{babel}
\usepackage[a4paper, top=1.6cm, bottom=1.6cm, left=1.4cm, right=1.4cm]{geometry}
\setlength{\headheight}{14.5pt}
\usepackage{amsmath, amssymb, amsfonts}
\usepackage{siunitx}
\sisetup{locale=FR, output-decimal-marker={,}, inter-unit-product=\ensuremath{{}\cdot{}}}
\usepackage{array, booktabs, multirow, tabularx}
\usepackage{xcolor}
\usepackage[most, breakable]{tcolorbox}
\usepackage{multicol}
\usepackage{enumitem}
\usepackage{mathtools}
\usepackage{lmodern}
\usepackage{microtype}
\usepackage{fancyhdr}
\usepackage{tikz}
\usetikzlibrary{positioning,arrows.meta}
\usepackage{lastpage}

% ── Couleurs par domaine ──────────────────────────────────────────────
\definecolor{colElec}{RGB}{26, 82, 148}       % Électrotechnique — bleu foncé
\definecolor{colElecLight}{RGB}{210, 225, 245} % fond bleu clair
\definecolor{colAuto}{RGB}{34, 120, 66}        % Automatique — vert
\definecolor{colMeca}{RGB}{185, 90, 20}        % Mécanique — orange
\definecolor{colTherm}{RGB}{160, 40, 40}       % Thermique — rouge
\definecolor{colQT}{RGB}{232, 244, 233}        % fond question type (vert très léger)
\definecolor{colJury}{RGB}{255, 243, 205}      % fond jury (jaune ambre)
\definecolor{colJuryFrame}{RGB}{200, 130, 0}   % bordure jury
\definecolor{colGrayLight}{RGB}{245, 245, 247} % fond gris neutre
\definecolor{colSeparator}{RGB}{180, 180, 180} % gris séparateur

% ── Style tcolorbox de section principale ────────────────────────────
\tcbset{
  ficheSection/.style={
    breakable,
    enhanced,
    colback=#1!7!white,
    colframe=#1,
    colbacktitle=#1,
    coltitle=white,
    fonttitle=\bfseries\normalsize,
    left=4pt, right=4pt, top=3pt, bottom=4pt,
    before skip=6pt, after skip=6pt,
  }
}

% ── Style encadré question type ──────────────────────────────────────
\tcbset{
  questiontype/.style={
    breakable,
    enhanced,
    colback=colQT,
    colframe=colElec!55!black,
    colbacktitle=colElec!55!black,
    coltitle=white,
    fonttitle=\bfseries\small,
    left=4pt, right=4pt, top=3pt, bottom=3pt,
    before skip=5pt, after skip=5pt,
  }
}

% ── Style encadré jury ───────────────────────────────────────────────
\tcbset{
  jury/.style={
    breakable,
    enhanced,
    colback=colJury,
    colframe=colJuryFrame,
    fonttitle=\bfseries\normalsize,
    left=4pt, right=4pt, top=3pt, bottom=3pt,
  }
}

% ── Mise en page ─────────────────────────────────────────────────────
\pagestyle{fancy}
\fancyhf{}
\renewcommand{\headrulewidth}{0.4pt}
\renewcommand{\footrulewidth}{0.3pt}

\setlist[itemize]{leftmargin=1.2em, itemsep=1pt, topsep=2pt, parsep=0pt}
\setlist[enumerate]{leftmargin=1.2em, itemsep=3pt, topsep=2pt, parsep=0pt}

\setlength{\parskip}{2pt}
\setlength{\columnsep}{10pt}

\def\CAPETBEGINFILE{\begin{document}}
\def\CAPETENDFILE{\end{document}}
\else
\def\CAPETBEGINFILE{}
\def\CAPETENDFILE{}
\fi

\fancyhead[L]{\textbf{\textcolor{colElec}{T07}} - Hacheurs (1, 2 et 4 quadrants)}
\fancyhead[R]{\small\textit{CAPET SII IE - Session 2026}}
\fancyfoot[C]{\small\thepage\,/\,\pageref{LastPage}}

\CAPETBEGINFILE

% -----------------------------------------------------------------------
%  BLOC EN-TeTE
% -----------------------------------------------------------------------

\begin{tcolorbox}[
  enhanced, breakable,
  colback=colElec,
  colframe=colElec,
  left=6pt, right=6pt, top=5pt, bottom=5pt,
  before skip=0pt, after skip=4pt
]
  {\color{white}\bfseries\large T07 \quad Hacheurs (1, 2 et 4 quadrants)}\\[2pt]
  {\color{colElec!30!white}\small
    \textbf{Domaine :} electrotechnique - electronique de puissance \quad
    \textbf{Presence :} \textbf{9/13} sujets (2013-2026) - theme \emph{tres frequent}
  }\\[1pt]
  {\color{colElec!20!white}\footnotesize
    Sous-themes : hacheur Buck (devolteur) - hacheur Boost (elevateur) - ondulation courant - hacheur 4 quadrants (pont en H) - entrelacement de bras - MPPT
  }
\end{tcolorbox}

% -----------------------------------------------------------------------
%  1. RAPPELS DE COURS
% -----------------------------------------------------------------------

\begin{tcolorbox}[
  ficheSection=colElec,
  title={1. Rappels de cours},
]

% -- 1.1 Hacheur Buck (devolteur) -------------------------------------
\textbf{1.1 Hacheur Buck - devolteur (serie)}

\smallskip
Principe : l'interrupteur commande est ferme pendant $t_{on} = \alpha T$ puis ouvert pendant $t_{off} = (1-\alpha)T$.
En conduction continue, la loi de la valeur moyenne donne :
\[
  \langle u_{sortie} \rangle = \alpha \cdot E \qquad \text{avec}\quad \alpha \in [0\,;\,1]
\]
Ondulation de courant dans l'inductance (pire cas - $\alpha = 0{,}5$) :
\[
  \Delta i_L = \frac{E \cdot \alpha \cdot (1-\alpha)}{L \cdot f}
  \qquad\Rightarrow\qquad
  L_{min} = \frac{E \cdot 0{,}5 \times 0{,}5}{\Delta i_{L,max} \cdot f} = \frac{E}{4 \cdot \Delta i_{L,max} \cdot f}
\]

\medskip
\hrule\vspace{4pt}

% -- 1.2 Hacheur Boost (elevateur) ------------------------------------
\textbf{1.2 Hacheur Boost - elevateur}

\smallskip
En regime permanent, $\langle v_L \rangle = 0$ sur une periode :
\[
  E_b \cdot \alpha + (E_b - V_{DC})(1-\alpha) = 0
  \quad\Rightarrow\quad
  \boxed{V_{DC} = \frac{E_b}{1-\alpha}} > E_b
\]
Phase $Q_1$ ferme ($t \in [0\,;\,\alpha T[$) : $v_{L} = E_b$ (courant croet).\\
Phase $Q_1$ ouvert ($t \in [\alpha T\,;\,T[$) : $v_{L} = E_b - V_{DC}$ (courant decroet, diode passe).

\medskip
\hrule\vspace{4pt}

% -- 1.3 Hacheur 4 quadrants (pont en H) ------------------------------
\textbf{1.3 Hacheur 4 quadrants - pont en H}

\smallskip
\begin{center}
\renewcommand{\arraystretch}{1.35}
\begin{tabular}{lll}
\toprule
\textbf{Configuration} & \textbf{Tension moyenne} & \textbf{Plage} \\
\midrule
Hacheur Buck (1Q) & $\langle u \rangle = \alpha \cdot E$ & $0 \leq \alpha \leq 1$ \\[4pt]
Pont en H bipolaire (4Q) & $\langle u \rangle = (2\alpha - 1) \cdot E$ & $\alpha = 0{,}5 \to u = 0$ \\[4pt]
Entrelacement $n$ bras & $\Delta i_{resultante} = \Delta i_L / n$ & Dephasage $T/n$ entre bras \\
\bottomrule
\end{tabular}
\end{center}

\smallskip
Conservation de la puissance (Boost ideal) : $E_b \cdot I_{E} = V_{DC} \cdot I_{DC}$.

\end{tcolorbox}

% -----------------------------------------------------------------------
%  2. FORMULAIRE SYNTHeTIQUE
% -----------------------------------------------------------------------

\begin{tcolorbox}[
  ficheSection=colElec,
  title={2. Formulaire synthetique},
]
\small
\begin{multicols}{2}
\renewcommand{\arraystretch}{1.45}
\begin{tabular}{@{}p{4.5cm}ll@{}}
\toprule
\textbf{Formule} & \textbf{Usage} & \textbf{U.} \\
\midrule
\multicolumn{3}{@{}l}{\footnotesize\textit{Hacheur Buck}} \\
$\langle u_s \rangle = \alpha \cdot E$ & Tension sortie & V \\
$\alpha = \langle u_s \rangle / E$ & Rapport cyclique & - \\
$\Delta i_L = E\alpha(1-\alpha)/(Lf)$ & Ondulation & A \\
$L_{min} = E/(4\Delta i_{max} f)$ & Inductance min & H \\
$t_{on} = \alpha T$ & Duree passant & s \\
\midrule
\multicolumn{3}{@{}l}{\footnotesize\textit{Hacheur Boost}} \\
$V_{DC} = E_b/(1-\alpha)$ & Tension sortie & V \\
$\langle v_L \rangle = 0$ & Regime perm. & - \\
$E_b I_E = V_{DC} I_{DC}$ & Bilan puiss. & W \\
\bottomrule
\end{tabular}

\columnbreak

\renewcommand{\arraystretch}{1.45}
\begin{tabular}{@{}p{4.5cm}ll@{}}
\toprule
\textbf{Formule} & \textbf{Usage} & \textbf{U.} \\
\midrule
\multicolumn{3}{@{}l}{\footnotesize\textit{Pont en H (4Q)}} \\
$\langle u \rangle = (2\alpha-1)E$ & Tens. moy. bip. & V \\
$\alpha > 0{,}5 \to u > 0$ & Avance & - \\
$\alpha < 0{,}5 \to u < 0$ & Revers. & - \\
\midrule
\multicolumn{3}{@{}l}{\footnotesize\textit{Entrelacement}} \\
$\Delta i_{res} = \Delta i_L / n$ & Ondulation res. & A \\
$n$ bras dephases de $T/n$ & Condition & - \\
\midrule
\multicolumn{3}{@{}l}{\footnotesize\textit{Donnees typiques}} \\
$f = 50\,\mathrm{kHz}$ & Freq. comm. & Hz \\
$\alpha_{max,\Delta i} = 0{,}5$ & Pire cas & - \\
\bottomrule
\end{tabular}
\end{multicols}
\normalsize
\end{tcolorbox}

% -----------------------------------------------------------------------
%  3. QUESTIONS TYPES
% -----------------------------------------------------------------------

\begin{tcolorbox}[
  ficheSection=colElec,
  title={3. Questions types - extraites des sujets reels},
]

% -- QT1 -------------------------------------------------------------

% Mini-procedure de resolution type copie concours
\begin{tcolorbox}[questiontype,title={Procedure type copie concours (5 etapes)}]
\begin{enumerate}[leftmargin=1.5em]
  \item \textbf{Qualifier le modele} : hypotheses, domaine de validite, grandeurs d'entree/sortie.
  \item \textbf{Ecrire les lois} : equations de base avec symboles et unites SI.
  \item \textbf{Isoler l'inconnue} : expression litterale, puis application numerique tracee.
  \item \textbf{Controler l'ordre de grandeur} : coherence physique et unite finale.
  \item \textbf{Conclure concours} : phrase courte orientee decision technique.
\end{enumerate}
\end{tcolorbox}

% Visuel de methode : chemin de resolution CAPET
\begin{center}
\begin{tikzpicture}[font=\small, node distance=7mm and 6mm, >=Stealth]
  \node[draw, rounded corners, fill=colElec!8, inner sep=4pt] (n1) {Notions};
  \node[draw, rounded corners, fill=colElec!8, right=of n1, inner sep=4pt] (n2) {Methode};
  \node[draw, rounded corners, fill=colElec!8, right=of n2, inner sep=4pt] (n3) {Application annale};
  \node[draw, rounded corners, fill=colElec!8, right=of n3, inner sep=4pt] (n4) {Verification jury};
  \draw[->, thick, colElec] (n1) -- (n2);
  \draw[->, thick, colElec] (n2) -- (n3);
  \draw[->, thick, colElec] (n3) -- (n4);
\end{tikzpicture}
\end{center}
\begin{tcolorbox}[questiontype,
  title={QT1 - Hacheur Buck - rapport cyclique et inductance minimale \hfill\normalfont\footnotesize E2 disciplinaire 2023 (robot Bakus)}
]
\textbf{enonce.}
Un convertisseur devolteur recharge une batterie Li-ion depuis une LiFePO4 ($E = 500\,\mathrm{V}$). Conduction continue, $\langle u_{sortie} \rangle \in [70\,;\,117{,}5]\,\mathrm{V}$, $\Delta i \leq 2\,\mathrm{A}$, $f = 50\,\mathrm{kHz}$.
\begin{enumerate}[label=(\alph*)]
  \item Exprimer $\langle u_{sortie} \rangle$ en fonction de $E$ et $\alpha$.
  \item Donner la plage de $\alpha$.
  \item Calculer $L_{min}$.
\end{enumerate}

\tcbline

\textbf{Corrige.}

\textbf{(a)} Pendant $[0;\alpha T[$ l'interrupteur conduit : $u_{sortie} = E$. Pendant $[\alpha T;T[$ la diode roue-libre conduit : $u_{sortie} = 0$.
\[
  \langle u_{sortie} \rangle = \alpha \cdot E
\]

\textbf{(b)} $\alpha_{min} = 70/500 = 0{,}14$ ; $\alpha_{max} = 117{,}5/500 = 0{,}235$.

\textbf{(c)} L'ondulation est maximale pour $\alpha = 0{,}5$ (pire cas) :
\[
  L_{min} = \frac{E}{4 \cdot \Delta i_{max} \cdot f} = \frac{500}{4 \times 2 \times 50\,000} = \frac{500}{400\,000} = \mathbf{1{,}25\,\mathrm{mH}}
\]
\end{tcolorbox}

\vspace{4pt}

% -- QT2 -------------------------------------------------------------

\begin{tcolorbox}[questiontype,
  title={QT2 - Hacheur Boost - derivation par lois de Kirchhoff \hfill\normalfont\footnotesize Rapport jury 2025 (Q27-Q32)}
]
\textbf{enonce.}
Pour un hacheur elevateur ($E_b$ source, $V_{DC}$ charge) avec interrupteur $Q_1$ et diode $D_1$ : ecrire les lois de Kirchhoff pendant les deux phases. En deduire $V_{DC} = E_b/(1-\alpha)$.

\tcbline

\textbf{Corrige.}

\textbf{Phase 1} ($Q_1$ ferme, $D_1$ bloquee) : $v_L = E_b > 0$ - courant $i_L$ croet.

\textbf{Phase 2} ($Q_1$ ouvert, $D_1$ passante) : $v_L = E_b - V_{DC} < 0$ - courant $i_L$ decroet.

\textbf{Regime permanent} : $\langle v_L \rangle = 0$ sur une periode :
\[
  E_b \cdot \alpha T + (E_b - V_{DC})(1-\alpha)T = 0
  \quad\Rightarrow\quad
  E_b = V_{DC}(1-\alpha)
  \quad\Rightarrow\quad
  \boxed{V_{DC} = \frac{E_b}{1-\alpha}}
\]
\textit{Le jury note : "Trop peu de candidats appliquent correctement les lois de Kirchhoff aux deux phases." - ecrire systematiquement $v_L$, $v_D$, $v_Q$ pour chaque etat.}
\end{tcolorbox}

\vspace{4pt}

% -- QT3 -------------------------------------------------------------

\begin{tcolorbox}[questiontype,
  title={QT3 - Entrelacement de 4 bras - ondulation resultante \hfill\normalfont\footnotesize E2 disciplinaire 2023 (question C7)}
]
\textbf{enonce.}
Le convertisseur Buck du robot Bakus est realise avec 4 bras entrelaces (dephasage $T/4$ entre bras), chaque bras ayant $L = 1{,}25\,\mathrm{mH}$. L'ondulation par bras est $\Delta i_L = 2\,\mathrm{A}$. Quelle est l'ondulation resultante du courant de charge - 
\tcbline

\textbf{Corrige.}

Avec $n = 4$ bras dephases de $T/4$, les ondulations se compensent partiellement.
L'ondulation resultante est divisee par $n$ :
\[
  \Delta i_{resultante} = \frac{\Delta i_L}{n} = \frac{2}{4} = \mathbf{0{,}5\,\mathrm{A}}
\]
\textbf{Avantage :} meme inductance $L$, ondulation 4 fois reduite. Ou bien : meme specification d'ondulation, inductance 4 fois reduite.\\
\textit{Condition necessaire : dephasage exactement $T/4$ entre chaque bras. Sans dephasage, les ondulations s'additionnent.}
\end{tcolorbox}

\end{tcolorbox}

% -----------------------------------------------------------------------
%  4. ATTENTES DU JURY
% -----------------------------------------------------------------------

\begin{tcolorbox}[
  jury,
  title={4. Attentes du jury - Erreurs recurrentes et points de vigilance},
]

\begin{enumerate}[label=\textbf{\arabic*.}, leftmargin=*, itemsep=5pt]

  \item \textbf{Confondre Buck et Boost.}
  Buck (devolteur) : $V_{out} = \alpha E < E$ (rapport cyclique $< 1$).
  Boost (elevateur) : $V_{out} = E/(1-\alpha) > E$. L'erreur classique est d'appliquer la formule Buck - un schema Boost.

  \item \textbf{Signe de $v_L$ en phase OFF du Boost.}
  $v_L = E_b - V_{DC} < 0$ : le courant \textbf{decroet}. Beaucoup de candidats obtiennent un courant croissant dans les deux phases en inversant le signe.

  \item \textbf{La condition $\langle v_L \rangle = 0$ est la cle.}
  Sans cette condition de regime permanent, il est impossible de deriver les formules de tension moyenne. L'enoncer explicitement avant tout calcul.

  \item \textbf{Ondulation maximale - $\alpha = 0{,}5$, pas - $\alpha$ nominal.}
  Calculer $L_{min}$ pour $\alpha = 0{,}5$ (pire cas), meme si le point de fonctionnement nominal est $\alpha = 0{,}14$.

  \item \textbf{Identifier la nature du convertisseur depuis le schema.}
  Le jury apprecie de nommer Buck/Boost/4Q et de justifier par le sens du flux d'energie (sens courant, rapport $V_{out}/V_{in}$).

  \item \textbf{Entrelacement : dephasage obligatoire.}
  L'ondulation est divisee par $n$ \textbf{seulement si} les bras sont dephases de $T/n$.

\end{enumerate}

\medskip
\begin{tcolorbox}[
  colback=colJuryFrame!12!colJury,
  colframe=colJuryFrame,
  left=4pt, right=4pt, top=2pt, bottom=2pt,
  fontupper=\small\itshape
]
Rapport jury 2025 : "Q27-Q32 : trop peu de candidats appliquent correctement les lois de Kirchhoff aux deux phases de commutation - ecrire systematiquement $v_L$, $v_D$, $v_Q$ pour les deux etats."\\
Rapport jury 2025 : "Q14 : CV1 = Hacheur, CV2 = Redresseur, CV3 = Onduleur - les fonctions explicites doivent etre nommees et justifiees."
\end{tcolorbox}

\end{tcolorbox}

% -----------------------------------------------------------------------
%  PIED - PReSENCE DANS LES SUJETS
% -----------------------------------------------------------------------

\begin{tcolorbox}[
  enhanced, breakable,
  colback=colGrayLight,
  colframe=colSeparator,
  left=4pt, right=4pt, top=3pt, bottom=3pt,
  before skip=4pt, after skip=0pt,
  fonttitle=\small\bfseries,
  title={Presence dans les sujets analyses (2013-2026)},
]
\begin{center}
\footnotesize\renewcommand{\arraystretch}{1.15}
\begin{tabular}{lccccccccccccc}
\toprule
\textbf{Annee} & 13 & 14 & 15 & 16 & 17 & 18 & 19 & 20 & 21 & 22 & 23 & 24 & 25 \\
\midrule
E1 (commune)   & - & - & $\checkmark$ & - & $\checkmark$ & - & - & $\checkmark$ & - & $\checkmark$ & - & - & $\checkmark$ \\
E2 discipl. IE & $\checkmark$ & - & - & $\checkmark$ & - & $\checkmark$ & $\checkmark$ & - & $\checkmark$ & $\checkmark$ & $\checkmark$ & $\checkmark$ & $\checkmark$ \\
\bottomrule
\end{tabular}
\end{center}
\smallskip
\footnotesize\textbf{Sous-themes frequents :}
hacheur Buck - rapport cyclique (8/13) - inductance minimale (5/13) - Boost / MPPT (4/13) - hacheur 4 quadrants / pont en H (3/13) - entrelacement de bras (2/13)

\medskip
\textbf{Ordres de grandeur (reperes rapides).}
\begin{itemize}[leftmargin=1.3em]
  \item Reseau BT FR : tension phase-neutre typique \SI{230}{\volt}, tension phase-phase typique \SI{400}{\volt}, frequence \SI{50}{\hertz}.
  \item Courants d'etude frequents en sujet : de quelques \si{\ampere} a quelques dizaines de \si{\ampere}.
  \item En conversion d'energie, un rendement attendu reste generalement entre 0{,}8 et 0{,}98 selon la technologie.
\end{itemize}

\medskip
\textbf{Synthese express (a refaire sans support).}
\begin{enumerate}[leftmargin=1.4em]
  \item Poser les hypotheses du modele et les conventions de signe.
  \item Ecrire 2 a 4 relations cles avec unites.
  \item Realiser un calcul numerique complet et verifier l'ordre de grandeur.
  \item Justifier le resultat par un argument technique de niveau CAPET.
\end{enumerate}

\end{tcolorbox}

\CAPETENDFILE









